%% ma: L12-B06-0040
%% lop: 12
%% chuong: 2
%% bai: 6
%% dang: 12.6.18
%% loai: TN4
%% muc: NB
%% dap_an: "D"
%% nguon: "LaTeX giáo viên gửi 21/10/2026 (M05, Vectơ và hệ trục tọa độ trong không gian, 107 câu), Phần I Câu 21"
%% kien_thuc: "Ba vectơ đồng phẳng, điều kiện đồng phẳng (không có trong SGK Toán 12 KNTT Bài 6; thuộc chương trình cũ)"
%% trang_thai: cach-ly
%% ly_do: "Vượt chương trình: SGK KNTT 12 Bài 6 không dạy khái niệm và điều kiện ba vectơ đồng phẳng. Ngoài ra phương án A (thiếu điều kiện a, b không cùng phương) và D (thêm giả thiết) đều chưa chặt nên đáp án gây tranh cãi. Gợi ý: bỏ câu, hoặc đổi thành câu về quy tắc hình hộp"
\begin{ex}
Điều kiện cần và đủ để ba vectơ $\vec{a}$, $\vec{b}$, $\vec{c}$ khác $\vec{0}$ đồng phẳng trong không gian là
\choice{Tồn tại duy nhất một cặp số thực $(m;n)$ sao cho $\vec{c}=m\vec{a}+n\vec{b}$}{Ba vectơ đôi một cùng phương}{Tổng ba vectơ bằng vectơ $\vec{0}$}{\True Tồn tại các số thực $m$, $n$ sao cho $\vec{c}=m\vec{a}+n\vec{b}$ (với $\vec{a}$, $\vec{b}$ không cùng phương)}
\loigiai{Khi $\vec{a}$, $\vec{b}$ không cùng phương: $\vec{a}$, $\vec{b}$, $\vec{c}$ đồng phẳng $\Leftrightarrow$ tồn tại $m$, $n$ để $\vec{c}=m\vec{a}+n\vec{b}$.}
\end{ex}
